\section{Related Work}
\label{sec:related_work}

We review related literature across three key areas: weight-space model merging, test-time adaptive merging, and neural network quantization and sharpness-aware robustness.

\subsection{Weight-Space Model Merging}
Weight-space model merging has emerged as an efficient, training-free paradigm to consolidate multiple independently fine-tuned expert models into a single multi-task network. Foundational work by \citet{model_soups} introduced model souping, showing that averaging weights of fine-tuned models on the same task improves generalization. To extend weight-space composition to different tasks, \citet{task_vectors} proposed Task Arithmetic, which adds fine-tuned task-specific ``task vectors'' to a shared pre-trained initialization. 

However, direct addition often suffers from interference and sign conflicts among parameters. To address this, \citet{ties_merging} proposed TIES-Merging, which resolves parameter conflicts by keeping only the top- $k$\% largest magnitude values, aligning sign directions, and averaging the remaining values. Similarly, \citet{zipmerge} and \citet{symerge} explored pruning and matching weights before composition. While successful, these methods rely on uniform, hand-tuned scaling factors, ignoring task conflicts at different layers.

\subsection{Test-Time Adaptive Model Merging}
To overcome the limitations of static merging coefficients, recent methods have proposed learning the blending factors dynamically. The pioneer in this space, \textbf{AdaMerging} \cite{adamerging}, formulates model merging as a test-time adaptation (TTA) problem. Operating entirely without the original training data, AdaMerging optimizes task-wise or layer-wise coefficients on a small calibration stream of unlabeled test queries by minimizing an unsupervised surrogate objective, specifically the softmax entropy of the merged model's predictions. This test-time adaptation loop mirrors classic single-model TTA methods such as Tent \cite{tent} and SHOT \cite{shot}. 

Subsequent extensions, such as AdaMerging++ and related benchmark methods in FusionBench \cite{fusionbench}, have shown that layer-wise coefficients significantly outperform task-wise coefficients by providing fine-grained, localized control over parameter composition. However, unregularized entropy minimization at test time is prone to severe transductive overfitting on small calibration streams, leading to sharp local minima that are highly fragile under weight perturbations.

\subsection{Quantization and Sharpness-Aware Robustness}
Post-training quantization (PTQ) is a critical technique for deploying modern deep neural networks on resource-constrained edge hardware, compressing parameters to low-precision formats such as INT8 or INT4 \cite{gholami2022survey}. While PTQ significantly reduces memory and latency, the rounding operations inject high-dimensional quantization noise into the parameters, which often causes catastrophic performance collapse in sensitive models.

A well-established line of research connects quantization robustness to the geometry of the loss landscape, specifically the flatness of the local minimum. \citet{sam} proposed Sharpness-Aware Minimization (SAM), which explicitly minimizes loss sharpness during training to improve generalization. \citet{asymmetric_flatness} and \citet{quantization_flatness} demonstrated that models converged to flat minima exhibit superior robustness to PTQ noise because the quadratic loss increase under weight perturbations is bounded by the Hessian's eigenvalues. While prior work has explored flatness for training individual networks, HessMerge is the first to establish and leverage this second-order geometric relationship for weight-space model merging and coefficient optimization.
